\Chapter{78}

\Verse{1}
{
    \zA{话说两个尼姑领了芳官等去后，王夫人便往贾母处来。}
}
{
    \eA{[English source not present in the checked-in Bencraft/Joly files.]}
}
{
}

\Verse{2}
{
    \zA{见贾母喜欢，便趁便回 道：“宝玉屋里有个晴雯，那个丫头也大了，而且一年之间病不离身。}
}
{
    \eA{[English source not present in the checked-in Bencraft/Joly files.]}
}
{
}

\Verse{3}
{
    \zA{我常见他比 别人分外淘气，也懒；前日又病倒了十几天，叫大夫瞧，说是女儿痨，所以我就赶 着叫他下去了。}
}
{
    \eA{[English source not present in the checked-in Bencraft/Joly files.]}
}
{
}

\Verse{4}
{
    \zA{若养好了，也不用叫他进来，就赏他家配人去也罢了。}
}
{
    \eA{[English source not present in the checked-in Bencraft/Joly files.]}
}
{
}

\Verse{5}
{
    \zA{再那几个学 戏的女孩子，我也做主放了：一则他们都会戏，口里没轻没重，只会混说，女孩儿 们听了，如何使得？}
}
{
    \eA{[English source not present in the checked-in Bencraft/Joly files.]}
}
{
}

\Verse{6}
{
    \zA{二则他们唱会子戏，白放了他们，也是应该的。}
}
{
    \eA{[English source not present in the checked-in Bencraft/Joly files.]}
}
{
}

\Verse{7}
{
    \zA{况丫头们也太 多，若说不够使，再挑上几个来，也是一样。”}
}
{
    \eA{[English source not present in the checked-in Bencraft/Joly files.]}
}
{
}

\Verse{8}
{
    \zA{贾母听了点头道：“这是正理，我 也正想着如此。}
}
{
    \eA{[English source not present in the checked-in Bencraft/Joly files.]}
}
{
}

\Verse{9}
{
    \zA{但晴雯这丫头，我看他甚好，言谈针线都不及他，将来还可以给宝 玉使唤的，谁知变了。”}
}
{
    \eA{[English source not present in the checked-in Bencraft/Joly files.]}
}
{
}

\Verse{10}
{
    \zA{王夫人笑道：“老太太挑中的人原不错，只是他命里没造 化，所以得了这个病。}
}
{
    \eA{[English source not present in the checked-in Bencraft/Joly files.]}
}
{
}

\Verse{11}
{
    \zA{俗语又说：‘女大十八变。’}
}
{
    \eA{[English source not present in the checked-in Bencraft/Joly files.]}
}
{
}

\Verse{12}
{
    \zA{况且有本事的人，未免就有些 调歪，老太太还有什么不曾经历过的？}
}
{
    \eA{[English source not present in the checked-in Bencraft/Joly files.]}
}
{
}

\Verse{13}
{
    \zA{三年前我也就留心这件事，先只取中了他。}
}
{
    \eA{[English source not present in the checked-in Bencraft/Joly files.]}
}
{
}

\Verse{14}
{
    \zA{我留心看了去，他色色比人强，只是不大沉重。}
}
{
    \eA{[English source not present in the checked-in Bencraft/Joly files.]}
}
{
}

\Verse{15}
{
    \zA{知大体，莫若袭人第一。}
}
{
    \eA{[English source not present in the checked-in Bencraft/Joly files.]}
}
{
}

\Verse{16}
{
    \zA{虽说贤妻 美妾，也要性情和顺，举止？}
}
{
    \eA{[English source not present in the checked-in Bencraft/Joly files.]}
}
{
}

\Verse{17}
{
    \zA{重的更好些。}
}
{
    \eA{[English source not present in the checked-in Bencraft/Joly files.]}
}
{
}

\Verse{18}
{
    \zA{袭人的模样虽比晴雯次一等，然放在房 里也算是一二等的。}
}
{
    \eA{[English source not present in the checked-in Bencraft/Joly files.]}
}
{
}

\Verse{19}
{
    \zA{况且行事大方，心地老实，这几年从未同着宝玉淘气。}
}
{
    \eA{[English source not present in the checked-in Bencraft/Joly files.]}
}
{
}

\Verse{20}
{
    \zA{凡宝玉 十分胡闹的事，他只有死劝的。}
}
{
    \eA{[English source not present in the checked-in Bencraft/Joly files.]}
}
{
}

\Verse{21}
{
    \zA{因此，品择了二年，一点不错了，我悄悄的把他丫 头的月钱止住，我的月分银子里批出二两银子来给他，不过使他自己知道，越发小 心效好之意。}
}
{
    \eA{[English source not present in the checked-in Bencraft/Joly files.]}
}
{
}

\Verse{22}
{
    \zA{且没有明说，一则宝玉年纪尚小，老爷知道了，又恐就耽误了书；二 则宝玉自以为自己跟前的人，不敢劝他说他，反倒纵性起来。}
}
{
    \eA{[English source not present in the checked-in Bencraft/Joly files.]}
}
{
}

\Verse{23}
{
    \zA{所以直到今日，才回 明老太太。”}
}
{
    \eA{[English source not present in the checked-in Bencraft/Joly files.]}
}
{
}

\Verse{24}
{
    \zA{贾母听了，笑道：“原来这样，如此更好了。}
}
{
    \eA{[English source not present in the checked-in Bencraft/Joly files.]}
}
{
}

\Verse{25}
{
    \zA{袭人本来从小儿不言不 语，我只说是‘没嘴的葫芦’。}
}
{
    \eA{[English source not present in the checked-in Bencraft/Joly files.]}
}
{
}

\Verse{26}
{
    \zA{既是你深知，岂有大错误的？”}
}
{
    \eA{[English source not present in the checked-in Bencraft/Joly files.]}
}
{
}

\Verse{27}
{
    \zA{王夫人又回今日贾 政如何夸奖，如何带他们逛去。}
}
{
    \eA{[English source not present in the checked-in Bencraft/Joly files.]}
}
{
}

\Verse{28}
{
    \zA{贾母听了，更加喜悦。}
}
{
    \eA{[English source not present in the checked-in Bencraft/Joly files.]}
}
{
}

\Verse{29}
{
    \zA{一时，只见迎春妆扮了前来告辞过去。}
}
{
    \eA{[English source not present in the checked-in Bencraft/Joly files.]}
}
{
}

\Verse{30}
{
    \zA{凤姐也来请早安，伺候早饭。}
}
{
    \eA{[English source not present in the checked-in Bencraft/Joly files.]}
}
{
}

\Verse{31}
{
    \zA{又说笑一 回，贾母歇晌，王夫人便唤了凤姐，问他丸药可曾配来。}
}
{
    \eA{[English source not present in the checked-in Bencraft/Joly files.]}
}
{
}

\Verse{32}
{
    \zA{凤姐道：“还不曾呢，如 今还是吃汤药。}
}
{
    \eA{[English source not present in the checked-in Bencraft/Joly files.]}
}
{
}

\Verse{33}
{
    \zA{太太只管放心，我已大好了。”}
}
{
    \eA{[English source not present in the checked-in Bencraft/Joly files.]}
}
{
}

\Verse{34}
{
    \zA{王夫人见他精神复初，也就信了， 因告诉撵逐晴雯等事。}
}
{
    \eA{[English source not present in the checked-in Bencraft/Joly files.]}
}
{
}

\Verse{35}
{
    \zA{又说：“宝丫头怎么私自回家去了？}
}
{
    \eA{[English source not present in the checked-in Bencraft/Joly files.]}
}
{
}

\Verse{36}
{
    \zA{你们都不知道？}
}
{
    \eA{[English source not present in the checked-in Bencraft/Joly files.]}
}
{
}

\Verse{37}
{
    \zA{我前儿顺 路都查了一查。}
}
{
    \eA{[English source not present in the checked-in Bencraft/Joly files.]}
}
{
}

\Verse{38}
{
    \zA{谁知兰小子的这一个新进来的奶子，也十分的妖调，也不喜欢他。}
}
{
    \eA{[English source not present in the checked-in Bencraft/Joly files.]}
}
{
}

\Verse{39}
{
    \zA{我说给你大嫂子了：好不好，叫他各自去罢。}
}
{
    \eA{[English source not present in the checked-in Bencraft/Joly files.]}
}
{
}

\Verse{40}
{
    \zA{我因问你大嫂子：‘宝丫头出去，难 道你们不知道吗？’}
}
{
    \eA{[English source not present in the checked-in Bencraft/Joly files.]}
}
{
}

\Verse{41}
{
    \zA{他说是告诉了他了，不两三日，等姨妈病好了就进来。}
}
{
    \eA{[English source not present in the checked-in Bencraft/Joly files.]}
}
{
}

\Verse{42}
{
    \zA{姨妈究 竟没什么大病，不过咳嗽腰疼，年年是如此的。}
}
{
    \eA{[English source not present in the checked-in Bencraft/Joly files.]}
}
{
}

\Verse{43}
{
    \zA{他这去的必有原故，不是有人得罪 了他了？}
}
{
    \eA{[English source not present in the checked-in Bencraft/Joly files.]}
}
{
}

\Verse{44}
{
    \zA{那孩子心重，亲戚们住一场，别得罪了人，反不好了。”}
}
{
    \eA{[English source not present in the checked-in Bencraft/Joly files.]}
}
{
}

\Verse{45}
{
    \zA{凤姐笑道：“谁 可好好的得罪着他？”}
}
{
    \eA{[English source not present in the checked-in Bencraft/Joly files.]}
}
{
}

\Verse{46}
{
    \zA{王夫人道：“别是宝玉有嘴无心，从来没个忌讳，高了兴信 嘴胡说也是有的。”}
}
{
    \eA{[English source not present in the checked-in Bencraft/Joly files.]}
}
{
}

\Verse{47}
{
    \zA{凤姐笑道：“这可是太太过于操心了。}
}
{
    \eA{[English source not present in the checked-in Bencraft/Joly files.]}
}
{
}

\Verse{48}
{
    \zA{若说他出去干正经事， 说正经话去，却像傻子；若只叫他进来，在这些姊妹跟前，以至于大小的丫头跟前， 最有尽让，又恐怕得罪了人，那是再不得有人恼他的。}
}
{
    \eA{[English source not present in the checked-in Bencraft/Joly files.]}
}
{
}

\Verse{49}
{
    \zA{我想薛妹妹此去必是为前夜 搜检众丫头的原故，他自然为信不及园里的人，他又是亲戚，现也有丫头老婆在内， 我们又不好去搜检。}
}
{
    \eA{[English source not present in the checked-in Bencraft/Joly files.]}
}
{
}

\Verse{50}
{
    \zA{他恐我们疑他，所以多了这个心，自己回避了。}
}
{
    \eA{[English source not present in the checked-in Bencraft/Joly files.]}
}
{
}

\Verse{51}
{
    \zA{也是应该避嫌 疑的。”}
}
{
    \eA{[English source not present in the checked-in Bencraft/Joly files.]}
}
{
}

\Verse{52}
{
    \zA{王夫人听了这话不错，自己遂低头一想，便命人去请了宝钗来，分晰前日 的事，以解他的疑心，又仍命他进来照旧居住。}
}
{
    \eA{[English source not present in the checked-in Bencraft/Joly files.]}
}
{
}

\Verse{53}
{
    \zA{宝钗陪笑道：“我原要早出去的， 因姨妈有许多大事，所以不便来说。}
}
{
    \eA{[English source not present in the checked-in Bencraft/Joly files.]}
}
{
}

\Verse{54}
{
    \zA{可巧前日妈妈又不好了，家里两个靠得的女人 又病，所以我趁便去了。}
}
{
    \eA{[English source not present in the checked-in Bencraft/Joly files.]}
}
{
}

\Verse{55}
{
    \zA{姨妈今日既已知道了，我正好回明，就从今日辞了，好搬 东西。”}
}
{
    \eA{[English source not present in the checked-in Bencraft/Joly files.]}
}
{
}

\Verse{56}
{
    \zA{王夫人凤姐都笑道：“你太固执了。}
}
{
    \eA{[English source not present in the checked-in Bencraft/Joly files.]}
}
{
}

\Verse{57}
{
    \zA{正经再搬进来为是，休为没要紧的事 反疏远了亲戚。”}
}
{
    \eA{[English source not present in the checked-in Bencraft/Joly files.]}
}
{
}

\Verse{58}
{
    \zA{宝钗笑道：“这话说的太重了，并没为什么事要出去。}
}
{
    \eA{[English source not present in the checked-in Bencraft/Joly files.]}
}
{
}

\Verse{59}
{
    \zA{我为的是 妈妈近来神思比先大减，而且夜晚没有得靠的人，统共只我一个人；二则如今我哥
        哥眼看娶嫂子，多少针线活计，并家里一切动用器皿，尚有未齐备的，我也须得帮 着妈妈去料理料理。}
}
{
    \eA{[English source not present in the checked-in Bencraft/Joly files.]}
}
{
}

\Verse{60}
{
    \zA{姨妈和凤姐姐都知道我们家的事，不是我撒谎。}
}
{
    \eA{[English source not present in the checked-in Bencraft/Joly files.]}
}
{
}

\Verse{61}
{
    \zA{再者，自我在 园里，东南上小角门子就常开着，原是为我走的，保不住出入的人图省走路，也从 那里走。}
}
{
    \eA{[English source not present in the checked-in Bencraft/Joly files.]}
}
{
}

\Verse{62}
{
    \zA{又没个人盘查，设若从那里弄出事来，岂不两碍？}
}
{
    \eA{[English source not present in the checked-in Bencraft/Joly files.]}
}
{
}

\Verse{63}
{
    \zA{而且我进园里来睡，原 不是什么大事。}
}
{
    \eA{[English source not present in the checked-in Bencraft/Joly files.]}
}
{
}

\Verse{64}
{
    \zA{因前几年年纪都小，且家里没事，在外头不如进来，姊妹们在一处 玩笑作针线，都比在外头一人闷坐好些。}
}
{
    \eA{[English source not present in the checked-in Bencraft/Joly files.]}
}
{
}

\Verse{65}
{
    \zA{如今彼此都大了，况姨娘这边历年皆遇不 遂心之事，所以那园子里，倘有一时照顾不到的，皆有关系。}
}
{
    \eA{[English source not present in the checked-in Bencraft/Joly files.]}
}
{
}

\Verse{66}
{
    \zA{惟有少几个人，就可 以少操些心了。}
}
{
    \eA{[English source not present in the checked-in Bencraft/Joly files.]}
}
{
}

\Verse{67}
{
    \zA{所以今日不但我决意辞去，此外还要劝姨娘：如今该减省的就减省 些，也不为失了大家的体统。}
}
{
    \eA{[English source not present in the checked-in Bencraft/Joly files.]}
}
{
}

\Verse{68}
{
    \zA{据我看，园里的这一项费用也竟可以免的，说不得当 日的话。}
}
{
    \eA{[English source not present in the checked-in Bencraft/Joly files.]}
}
{
}

\Verse{69}
{
    \zA{姨娘深知我家的，难道我家当日也是这样零落不成？”}
}
{
    \eA{[English source not present in the checked-in Bencraft/Joly files.]}
}
{
}

\Verse{70}
{
    \zA{凤姐听了这篇话， 便向王夫人笑道：“这话依我竟不必强他。”}
}
{
    \eA{[English source not present in the checked-in Bencraft/Joly files.]}
}
{
}

\Verse{71}
{
    \zA{王夫人点头道：“我也无可回答，只 好随你的便罢了。”}
}
{
    \eA{[English source not present in the checked-in Bencraft/Joly files.]}
}
{
}

\Verse{72}
{
    \zA{说话之间，只见宝玉已回来了，因说：“老爷还未散，恐天黑了，所以先叫我 们回来了。”}
}
{
    \eA{[English source not present in the checked-in Bencraft/Joly files.]}
}
{
}

\Verse{73}
{
    \zA{王夫人忙问：“今日可丢了丑了没有？”}
}
{
    \eA{[English source not present in the checked-in Bencraft/Joly files.]}
}
{
}

\Verse{74}
{
    \zA{宝玉笑道：“不但不丢丑， 拐了许多东西来。”}
}
{
    \eA{[English source not present in the checked-in Bencraft/Joly files.]}
}
{
}

\Verse{75}
{
    \zA{接着就有老婆子们从二门上小厮手内接进东西来。}
}
{
    \eA{[English source not present in the checked-in Bencraft/Joly files.]}
}
{
}

\Verse{76}
{
    \zA{王夫人一看 时，只见扇子三把，扇坠三个，笔墨共六匣，香珠三串，玉绦环三个。}
}
{
    \eA{[English source not present in the checked-in Bencraft/Joly files.]}
}
{
}

\Verse{77}
{
    \zA{宝玉说道： “这是梅翰林送的，那是杨侍郎送的，这是李员外送的：每人一分。”}
}
{
    \eA{[English source not present in the checked-in Bencraft/Joly files.]}
}
{
}

\Verse{78}
{
    \zA{说着，又向 怀中取出一个檀香小护身佛来，说：“这是庆国公单给我的。”}
}
{
    \eA{[English source not present in the checked-in Bencraft/Joly files.]}
}
{
}

\Verse{79}
{
    \zA{王夫人又问在席何 人，做何诗词。}
}
{
    \eA{[English source not present in the checked-in Bencraft/Joly files.]}
}
{
}

\Verse{80}
{
    \zA{说毕，只将宝玉一分令人拿着，同宝玉、环、兰前来见贾母。}
}
{
    \eA{[English source not present in the checked-in Bencraft/Joly files.]}
}
{
}

\Verse{81}
{
    \zA{贾母 看了，喜欢不尽，不免又问些话。}
}
{
    \eA{[English source not present in the checked-in Bencraft/Joly files.]}
}
{
}

\Verse{82}
{
    \zA{无奈宝玉一心记着晴雯，答应完了，便说：“骑 马颠了，骨头疼。”}
}
{
    \eA{[English source not present in the checked-in Bencraft/Joly files.]}
}
{
}

\Verse{83}
{
    \zA{贾母便说：“快回房去，换了衣服，疏散疏散就好了，不许睡。”}
}
{
    \eA{[English source not present in the checked-in Bencraft/Joly files.]}
}
{
}

\Verse{84}
{
    \zA{宝玉听了，便忙进园来。}
}
{
    \eA{[English source not present in the checked-in Bencraft/Joly files.]}
}
{
}

\Verse{85}
{
    \zA{当下麝月秋纹已带了两个丫头来等候。}
}
{
    \eA{[English source not present in the checked-in Bencraft/Joly files.]}
}
{
}

\Verse{86}
{
    \zA{见宝玉辞了贾母出来，秋纹便将墨笔等 物拿着，随宝玉进园来。}
}
{
    \eA{[English source not present in the checked-in Bencraft/Joly files.]}
}
{
}

\Verse{87}
{
    \zA{宝玉满口里说：“好热。”}
}
{
    \eA{[English source not present in the checked-in Bencraft/Joly files.]}
}
{
}

\Verse{88}
{
    \zA{一壁走一面便摘冠解带，将外 面的大衣服都脱下来麝月拿着，只穿着一件松花绫子夹袄，襟内露出血点般大红裤 子来。}
}
{
    \eA{[English source not present in the checked-in Bencraft/Joly files.]}
}
{
}

\Verse{89}
{
    \zA{秋纹见这条红裤是晴雯针线，因叹道：“真是‘物在人亡’了！”}
}
{
    \eA{[English source not present in the checked-in Bencraft/Joly files.]}
}
{
}

\Verse{90}
{
    \zA{麝月将秋 纹拉了一把，笑道：“这裤子配着松花色袄儿、石青靴子，越显出靛青的头，雪白 的脸来了。”}
}
{
    \eA{[English source not present in the checked-in Bencraft/Joly files.]}
}
{
}

\Verse{91}
{
    \zA{宝玉在前，只装没听见，又走了两步便止步道：“我要走一走，这怎 么好？”}
}
{
    \eA{[English source not present in the checked-in Bencraft/Joly files.]}
}
{
}

\Verse{92}
{
    \zA{麝月道：“大白日里还怕什么，还怕丢了你不成？”}
}
{
    \eA{[English source not present in the checked-in Bencraft/Joly files.]}
}
{
}

\Verse{93}
{
    \zA{因命两个小丫头跟着， “我们送了这些东西去再来。”}
}
{
    \eA{[English source not present in the checked-in Bencraft/Joly files.]}
}
{
}

\Verse{94}
{
    \zA{宝玉道：“好姐姐，等一等我再去。”}
}
{
    \eA{[English source not present in the checked-in Bencraft/Joly files.]}
}
{
}

\Verse{95}
{
    \zA{麝月道：“我 们去了就来。}
}
{
    \eA{[English source not present in the checked-in Bencraft/Joly files.]}
}
{
}

\Verse{96}
{
    \zA{两个人手里都有东西，倒像摆执事的，一个捧着文房四宝，一个捧着 冠袍带履，成个什么样子。”}
}
{
    \eA{[English source not present in the checked-in Bencraft/Joly files.]}
}
{
}

\Verse{97}
{
    \zA{宝玉听了，正中心怀，便让他二人去了。}
}
{
    \eA{[English source not present in the checked-in Bencraft/Joly files.]}
}
{
}

\Verse{98}
{
    \zA{他便带了两个小丫头到一块山子石后 头，悄问他二人道：“自我去了，你袭人姐姐打发人去瞧晴雯姐姐没有？”}
}
{
    \eA{[English source not present in the checked-in Bencraft/Joly files.]}
}
{
}

\Verse{99}
{
    \zA{这一个 答道：“打发宋妈瞧去了。”}
}
{
    \eA{[English source not present in the checked-in Bencraft/Joly files.]}
}
{
}

\Verse{100}
{
    \zA{宝玉道：“回来说什么？”}
}
{
    \eA{[English source not present in the checked-in Bencraft/Joly files.]}
}
{
}

\Verse{101}
{
    \zA{小丫头道：“回来说：晴 雯姐姐直着脖子叫了一夜，今日早起，就闭了眼住了口，世事不知，只有倒气的分 儿了。”}
}
{
    \eA{[English source not present in the checked-in Bencraft/Joly files.]}
}
{
}

\Verse{102}
{
    \zA{宝玉忙道：“一夜叫的是谁？”}
}
{
    \eA{[English source not present in the checked-in Bencraft/Joly files.]}
}
{
}

\Verse{103}
{
    \zA{小丫头道：“一夜叫的是娘。”}
}
{
    \eA{[English source not present in the checked-in Bencraft/Joly files.]}
}
{
}

\Verse{104}
{
    \zA{宝玉拭泪 道：“还叫谁？”}
}
{
    \eA{[English source not present in the checked-in Bencraft/Joly files.]}
}
{
}

\Verse{105}
{
    \zA{小丫头说：“没有听见叫别人了。”}
}
{
    \eA{[English source not present in the checked-in Bencraft/Joly files.]}
}
{
}

\Verse{106}
{
    \zA{宝玉道：“你糊涂。}
}
{
    \eA{[English source not present in the checked-in Bencraft/Joly files.]}
}
{
}

\Verse{107}
{
    \zA{想必没 有听真。”}
}
{
    \eA{[English source not present in the checked-in Bencraft/Joly files.]}
}
{
}

\Verse{108}
{
    \zA{旁边那一个小丫头最伶俐，听宝玉如此说，便上来说：“真个他糊涂！”}
}
{
    \eA{[English source not present in the checked-in Bencraft/Joly files.]}
}
{
}

\Verse{109}
{
    \zA{又向宝玉说：“不但我听的真切，我还亲自偷着看去来着。”}
}
{
    \eA{[English source not present in the checked-in Bencraft/Joly files.]}
}
{
}

\Verse{110}
{
    \zA{宝玉听说，忙问：“怎 么又亲自看去？”}
}
{
    \eA{[English source not present in the checked-in Bencraft/Joly files.]}
}
{
}

\Verse{111}
{
    \zA{小丫头道：“我想，晴雯姐姐素日和别人不同，待我们极好。}
}
{
    \eA{[English source not present in the checked-in Bencraft/Joly files.]}
}
{
}

\Verse{112}
{
    \zA{如 今他虽受了委屈出去，我们不能别的法子救他，只亲去瞧瞧，也不枉素日疼我们一 场。}
}
{
    \eA{[English source not present in the checked-in Bencraft/Joly files.]}
}
{
}

\Verse{113}
{
    \zA{就是人知道了，回了太太，打我们一顿，也是愿受的。}
}
{
    \eA{[English source not present in the checked-in Bencraft/Joly files.]}
}
{
}

\Verse{114}
{
    \zA{所以我拚着一顿打，偷 着出去瞧了一瞧。}
}
{
    \eA{[English source not present in the checked-in Bencraft/Joly files.]}
}
{
}

\Verse{115}
{
    \zA{谁知他平生为人聪明，至死不变，见我去了，便睁开眼拉我的手 问：‘宝玉那里去了？’}
}
{
    \eA{[English source not present in the checked-in Bencraft/Joly files.]}
}
{
}

\Verse{116}
{
    \zA{我告诉他了。}
}
{
    \eA{[English source not present in the checked-in Bencraft/Joly files.]}
}
{
}

\Verse{117}
{
    \zA{他叹了一口气，说：‘不能见了！’}
}
{
    \eA{[English source not present in the checked-in Bencraft/Joly files.]}
}
{
}

\Verse{118}
{
    \zA{我就说： ‘姐姐何不等一等他回来见一面？’}
}
{
    \eA{[English source not present in the checked-in Bencraft/Joly files.]}
}
{
}

\Verse{119}
{
    \zA{他就笑道：‘你们不知道，我不是死：如今天 上少了一个花神，玉皇爷叫我去管花儿。}
}
{
    \eA{[English source not present in the checked-in Bencraft/Joly files.]}
}
{
}

\Verse{120}
{
    \zA{我如今在未正二刻就上任去了，宝玉须得 未正三刻才到家，只少一刻儿的工夫，不能见面。}
}
{
    \eA{[English source not present in the checked-in Bencraft/Joly files.]}
}
{
}

\Verse{121}
{
    \zA{世上凡有该死的人，阎王勾取了 去，是差些个小鬼来拿他的魂儿。}
}
{
    \eA{[English source not present in the checked-in Bencraft/Joly files.]}
}
{
}

\Verse{122}
{
    \zA{要迟延一时半刻，不过烧些纸浇些浆饭，那鬼只 顾抢钱去了，该死的人就可挨磨些工夫。}
}
{
    \eA{[English source not present in the checked-in Bencraft/Joly files.]}
}
{
}

\Verse{123}
{
    \zA{我这如今是天上的神仙来请，那里捱得时 刻呢？’}
}
{
    \eA{[English source not present in the checked-in Bencraft/Joly files.]}
}
{
}

\Verse{124}
{
    \zA{我听了这话，竟不大信。}
}
{
    \eA{[English source not present in the checked-in Bencraft/Joly files.]}
}
{
}

\Verse{125}
{
    \zA{及进来到屋里，留神看时辰表，果然是未正二刻， 他咽了气；正三刻上，就有人来叫我们说你来了。”}
}
{
    \eA{[English source not present in the checked-in Bencraft/Joly files.]}
}
{
}

\Verse{126}
{
    \zA{宝玉忙道：“你不认得字，所 以不知道，这原是有的。}
}
{
    \eA{[English source not present in the checked-in Bencraft/Joly files.]}
}
{
}

\Verse{127}
{
    \zA{不但花有一花神，还有总花神。}
}
{
    \eA{[English source not present in the checked-in Bencraft/Joly files.]}
}
{
}

\Verse{128}
{
    \zA{但他不知做总花神去了， 还是单管一样花神？”}
}
{
    \eA{[English source not present in the checked-in Bencraft/Joly files.]}
}
{
}

\Verse{129}
{
    \zA{这丫头听了，一时诌不来。}
}
{
    \eA{[English source not present in the checked-in Bencraft/Joly files.]}
}
{
}

\Verse{130}
{
    \zA{恰好这是八月时节，园中池上芙 蓉正开，这丫头便见景生情，忙答道：“我已曾问他：‘是管什么花的神？}
}
{
    \eA{[English source not present in the checked-in Bencraft/Joly files.]}
}
{
}

\Verse{131}
{
    \zA{告诉我 们，日后也好供养的。’}
}
{
    \eA{[English source not present in the checked-in Bencraft/Joly files.]}
}
{
}

\Verse{132}
{
    \zA{他说：‘你只可告诉宝玉一人，除他之外，不可泄了天机。’}
}
{
    \eA{[English source not present in the checked-in Bencraft/Joly files.]}
}
{
}

\Verse{133}
{
    \zA{就告诉我说，他就是专管芙蓉花的。”}
}
{
    \eA{[English source not present in the checked-in Bencraft/Joly files.]}
}
{
}

\Verse{134}
{
    \zA{宝玉听了这话，不但不为怪，亦且去悲生喜，便回过头来，看着那芙蓉笑道： “此花也须得这样一个人去主管。}
}
{
    \eA{[English source not present in the checked-in Bencraft/Joly files.]}
}
{
}

\Verse{135}
{
    \zA{我就料定他那样的人必有一番事业!虽然超生苦 海，从此再不能相见了。”}
}
{
    \eA{[English source not present in the checked-in Bencraft/Joly files.]}
}
{
}

\Verse{136}
{
    \zA{免不得伤感思念；因又想：“虽然临终未见，如今且去 灵前一拜，也算尽这五六年的情意。”}
}
{
    \eA{[English source not present in the checked-in Bencraft/Joly files.]}
}
{
}

\Verse{137}
{
    \zA{想毕，忙至屋里，正值麝月秋纹找来。}
}
{
    \eA{[English source not present in the checked-in Bencraft/Joly files.]}
}
{
}

\Verse{138}
{
    \zA{宝玉 又自穿戴了，只说去看黛玉，遂一人出园，往前次看望之处来。}
}
{
    \eA{[English source not present in the checked-in Bencraft/Joly files.]}
}
{
}

\Verse{139}
{
    \zA{意为停柩在内，谁 知他哥嫂见他一咽气，便回了进去，希图早早些得几两发送例银。}
}
{
    \eA{[English source not present in the checked-in Bencraft/Joly files.]}
}
{
}

\Verse{140}
{
    \zA{王夫人闻知，便 命赏了十两银子，又命：“即刻送到外头焚化了罢。}
}
{
    \eA{[English source not present in the checked-in Bencraft/Joly files.]}
}
{
}

\Verse{141}
{
    \zA{女子痨死的，断不可留！”}
}
{
    \eA{[English source not present in the checked-in Bencraft/Joly files.]}
}
{
}

\Verse{142}
{
    \zA{他 哥嫂听了这话，一面得银，一面催人立刻入殓，抬往城外化人厂上去了。}
}
{
    \eA{[English source not present in the checked-in Bencraft/Joly files.]}
}
{
}

\Verse{143}
{
    \zA{剩的衣裳 簪环，约有三四百金之数，他哥嫂自收了，为后日之计。}
}
{
    \eA{[English source not present in the checked-in Bencraft/Joly files.]}
}
{
}

\Verse{144}
{
    \zA{二人将门锁上，一同送殡 去了。}
}
{
    \eA{[English source not present in the checked-in Bencraft/Joly files.]}
}
{
}

\Verse{145}
{
    \zA{宝玉走来扑了一个空，站了半天，并无别法，只得复身进入园中。}
}
{
    \eA{[English source not present in the checked-in Bencraft/Joly files.]}
}
{
}

\Verse{146}
{
    \zA{及回至房中， 甚觉无味，因顺路来找黛玉，不在房里。}
}
{
    \eA{[English source not present in the checked-in Bencraft/Joly files.]}
}
{
}

\Verse{147}
{
    \zA{问其何往，丫鬟们回说：“往宝姑娘那里 去了。”}
}
{
    \eA{[English source not present in the checked-in Bencraft/Joly files.]}
}
{
}

\Verse{148}
{
    \zA{宝玉又至蘅芜院中，只见寂静无人，房内搬出，空空落落，不觉吃一大惊， 才想起前日仿佛听见宝钗要搬出去，只因这两日工课忙就混忘了，这时看见如此，
        才知道果然搬出。}
}
{
    \eA{[English source not present in the checked-in Bencraft/Joly files.]}
}
{
}

\Verse{149}
{
    \zA{怔了半天，因转念一想：“不如还是和袭人厮混，再与黛玉相伴。}
}
{
    \eA{[English source not present in the checked-in Bencraft/Joly files.]}
}
{
}

\Verse{150}
{
    \zA{只这两三个人，只怕还是同死同归。”}
}
{
    \eA{[English source not present in the checked-in Bencraft/Joly files.]}
}
{
}

\Verse{151}
{
    \zA{想毕，仍往潇湘馆来。}
}
{
    \eA{[English source not present in the checked-in Bencraft/Joly files.]}
}
{
}

\Verse{152}
{
    \zA{偏黛玉还未回来。}
}
{
    \eA{[English source not present in the checked-in Bencraft/Joly files.]}
}
{
}

\Verse{153}
{
    \zA{正 在不知所之，忽见王夫人的丫头进来找他，说：“老爷回来了，找你呢。}
}
{
    \eA{[English source not present in the checked-in Bencraft/Joly files.]}
}
{
}

\Verse{154}
{
    \zA{又得了好 题目了。}
}
{
    \eA{[English source not present in the checked-in Bencraft/Joly files.]}
}
{
}

\Verse{155}
{
    \zA{快走，快走。”}
}
{
    \eA{[English source not present in the checked-in Bencraft/Joly files.]}
}
{
}

\Verse{156}
{
    \zA{宝玉听了，只得跟了出来。}
}
{
    \eA{[English source not present in the checked-in Bencraft/Joly files.]}
}
{
}

\Verse{157}
{
    \zA{到王夫人屋里，他父亲已出去 了，王夫人命人送宝玉至书房里。}
}
{
    \eA{[English source not present in the checked-in Bencraft/Joly files.]}
}
{
}

\Verse{158}
{
    \zA{彼时贾政正与众幕友们谈论寻书之胜。}
}
{
    \eA{[English source not present in the checked-in Bencraft/Joly files.]}
}
{
}

\Verse{159}
{
    \zA{又说：“临散时，忽谈及一事，最是千 古佳谈，‘风流隽逸，忠义感慨’，八字皆备。}
}
{
    \eA{[English source not present in the checked-in Bencraft/Joly files.]}
}
{
}

\Verse{160}
{
    \zA{倒是个好题目，大家要做一首挽词。”}
}
{
    \eA{[English source not present in the checked-in Bencraft/Joly files.]}
}
{
}

\Verse{161}
{
    \zA{众幕宾听了，都请教：“系何等妙事？”}
}
{
    \eA{[English source not present in the checked-in Bencraft/Joly files.]}
}
{
}

\Verse{162}
{
    \zA{贾政乃道：“当日曾有一位王爵，封曰恒 王，出镇青州。}
}
{
    \eA{[English source not present in the checked-in Bencraft/Joly files.]}
}
{
}

\Verse{163}
{
    \zA{这恒王最喜女色，且公馀好武，因选了许多美女，日习武事，令众 美女学习战攻斗伐之事。}
}
{
    \eA{[English source not present in the checked-in Bencraft/Joly files.]}
}
{
}

\Verse{164}
{
    \zA{内中有个姓林行四的，姿色既佳，且武艺更精，皆呼为林 四娘。}
}
{
    \eA{[English source not present in the checked-in Bencraft/Joly files.]}
}
{
}

\Verse{165}
{
    \zA{恒王最得意，遂超拔林四娘统辖诸姬，又呼为？？}
}
{
    \eA{[English source not present in the checked-in Bencraft/Joly files.]}
}
{
}

\Verse{166}
{
    \zA{将军。”}
}
{
    \eA{[English source not present in the checked-in Bencraft/Joly files.]}
}
{
}

\Verse{167}
{
    \zA{众清客都称：“妙 极神奇。}
}
{
    \eA{[English source not present in the checked-in Bencraft/Joly files.]}
}
{
}

\Verse{168}
{
    \zA{竟以‘？？’}
}
{
    \eA{[English source not present in the checked-in Bencraft/Joly files.]}
}
{
}

\Verse{169}
{
    \zA{下加‘将军’二字，反更觉妩媚风流，真绝世奇文也。}
}
{
    \eA{[English source not present in the checked-in Bencraft/Joly files.]}
}
{
}

\Verse{170}
{
    \zA{想这 恒王也是千古第一风流人物了。”}
}
{
    \eA{[English source not present in the checked-in Bencraft/Joly files.]}
}
{
}

\Verse{171}
{
    \zA{贾政笑道：“这话自然如此。}
}
{
    \eA{[English source not present in the checked-in Bencraft/Joly files.]}
}
{
}

\Verse{172}
{
    \zA{但更有可奇可叹之 事。”}
}
{
    \eA{[English source not present in the checked-in Bencraft/Joly files.]}
}
{
}

\Verse{173}
{
    \zA{众清客都惊问道：“不知底下有何等奇事？”}
}
{
    \eA{[English source not present in the checked-in Bencraft/Joly files.]}
}
{
}

\Verse{174}
{
    \zA{贾政道：“谁知次年，便有‘黄 巾’‘赤眉’一干流贼馀党复又乌合，抢掠山左一带。}
}
{
    \eA{[English source not present in the checked-in Bencraft/Joly files.]}
}
{
}

\Verse{175}
{
    \zA{恒王意为犬羊之辈，不足大 举，因轻骑进剿。}
}
{
    \eA{[English source not present in the checked-in Bencraft/Joly files.]}
}
{
}

\Verse{176}
{
    \zA{不意贼众诡谲，两战不胜，恒王遂被众贼所戮。}
}
{
    \eA{[English source not present in the checked-in Bencraft/Joly files.]}
}
{
}

\Verse{177}
{
    \zA{于是青州城内文 武官员，各各皆谓：‘王尚不胜，你我何为？’}
}
{
    \eA{[English source not present in the checked-in Bencraft/Joly files.]}
}
{
}

\Verse{178}
{
    \zA{遂将有献城之举。}
}
{
    \eA{[English source not present in the checked-in Bencraft/Joly files.]}
}
{
}

\Verse{179}
{
    \zA{林四娘得闻凶信， 遂聚集众女将，发令说道：‘你我皆向蒙王恩，戴天履地，不能报其万一。}
}
{
    \eA{[English source not present in the checked-in Bencraft/Joly files.]}
}
{
}

\Verse{180}
{
    \zA{今王既 殒身国患，我意亦当殒身于下。}
}
{
    \eA{[English source not present in the checked-in Bencraft/Joly files.]}
}
{
}

\Verse{181}
{
    \zA{尔等有愿随着，即同我前往，不愿者亦早自散去。’}
}
{
    \eA{[English source not present in the checked-in Bencraft/Joly files.]}
}
{
}

\Verse{182}
{
    \zA{众女将听他这样，都一齐说：‘愿意！’}
}
{
    \eA{[English source not present in the checked-in Bencraft/Joly files.]}
}
{
}

\Verse{183}
{
    \zA{于是林四娘带领众人，连夜出城，直杀至 贼营。}
}
{
    \eA{[English source not present in the checked-in Bencraft/Joly files.]}
}
{
}

\Verse{184}
{
    \zA{里头众贼不防，也被斩杀了几个首贼。}
}
{
    \eA{[English source not present in the checked-in Bencraft/Joly files.]}
}
{
}

\Verse{185}
{
    \zA{后来大家见是不过几个女人，料不能 济事，遂回戈倒兵，奋力一阵，把林四娘等一个不曾留下，倒作成了这林四娘的一 片忠心之志。}
}
{
    \eA{[English source not present in the checked-in Bencraft/Joly files.]}
}
{
}

\Verse{186}
{
    \zA{后来报至都中，天子百官，无不叹息。}
}
{
    \eA{[English source not present in the checked-in Bencraft/Joly files.]}
}
{
}

\Verse{187}
{
    \zA{想其朝中自然又有人去剿灭， 天兵一到，化为乌有，不必深论。}
}
{
    \eA{[English source not present in the checked-in Bencraft/Joly files.]}
}
{
}

\Verse{188}
{
    \zA{只就林四娘一节，众位听了，可羡不可羡？”}
}
{
    \eA{[English source not present in the checked-in Bencraft/Joly files.]}
}
{
}

\Verse{189}
{
    \zA{众 幕友都叹道：“实在可羡可奇!实是个妙题，原该大家挽一挽才是。”}
}
{
    \eA{[English source not present in the checked-in Bencraft/Joly files.]}
}
{
}

\Verse{190}
{
    \zA{说着，早有 人取了笔砚，按贾政口中之言，稍加改易了几个字，便成了一篇短序，递给贾政看 了，贾政道：“不过如此。}
}
{
    \eA{[English source not present in the checked-in Bencraft/Joly files.]}
}
{
}

\Verse{191}
{
    \zA{他们那里已有原序。}
}
{
    \eA{[English source not present in the checked-in Bencraft/Joly files.]}
}
{
}

\Verse{192}
{
    \zA{昨日内又奉恩旨：着察核前代以来 应加褒奖而遗落未经奏请各项人等，无论僧、尼、乞丐、女妇人等，有一事可嘉， 即行汇送履历至礼部，备请恩奖。}
}
{
    \eA{[English source not present in the checked-in Bencraft/Joly files.]}
}
{
}

\Verse{193}
{
    \zA{所以他这原序也送往礼部去了。}
}
{
    \eA{[English source not present in the checked-in Bencraft/Joly files.]}
}
{
}

\Verse{194}
{
    \zA{大家听了这新闻， 所以都要做一首《？？}
}
{
    \eA{[English source not present in the checked-in Bencraft/Joly files.]}
}
{
}

\Verse{195}
{
    \zA{词》，以志其忠义。”}
}
{
    \eA{[English source not present in the checked-in Bencraft/Joly files.]}
}
{
}

\Verse{196}
{
    \zA{众人听了，都又笑道：“这原该如此。}
}
{
    \eA{[English source not present in the checked-in Bencraft/Joly files.]}
}
{
}

\Verse{197}
{
    \zA{只是更可羡者，本朝皆系千古未有之旷典，可谓‘圣朝无阙事’了。”}
}
{
    \eA{[English source not present in the checked-in Bencraft/Joly files.]}
}
{
}

\Verse{198}
{
    \zA{贾政点头道： “正是。”}
}
{
    \eA{[English source not present in the checked-in Bencraft/Joly files.]}
}
{
}

\Verse{199}
{
    \zA{说话间，宝玉、贾环、贾兰俱起身来看了题目。}
}
{
    \eA{[English source not present in the checked-in Bencraft/Joly files.]}
}
{
}

\Verse{200}
{
    \zA{贾政命他三人各吊一首，谁先 做成者赏，佳者额外加赏。}
}
{
    \eA{[English source not present in the checked-in Bencraft/Joly files.]}
}
{
}

\Verse{201}
{
    \zA{贾环贾兰二人近日当着许多人皆做过几首了，胆量愈壮。}
}
{
    \eA{[English source not present in the checked-in Bencraft/Joly files.]}
}
{
}

\Verse{202}
{
    \zA{今看了题目，遂自去思索。}
}
{
    \eA{[English source not present in the checked-in Bencraft/Joly files.]}
}
{
}

\Verse{203}
{
    \zA{一时贾兰先有了，贾环生恐落后，也就有了。}
}
{
    \eA{[English source not present in the checked-in Bencraft/Joly files.]}
}
{
}

\Verse{204}
{
    \zA{二人皆已 录出，宝玉尚自出神。}
}
{
    \eA{[English source not present in the checked-in Bencraft/Joly files.]}
}
{
}

\Verse{205}
{
    \zA{贾政与众人且看他二人的二首。}
}
{
    \eA{[English source not present in the checked-in Bencraft/Joly files.]}
}
{
}

\Verse{206}
{
    \zA{贾兰的是一首七言绝句，写道是： ？？？？}
}
{
    \eA{[English source not present in the checked-in Bencraft/Joly files.]}
}
{
}

\Verse{207}
{
    \zA{将军林四娘，玉为肌骨铁为肠。}
}
{
    \eA{[English source not present in the checked-in Bencraft/Joly files.]}
}
{
}

\Verse{208}
{
    \zA{捐躯自报恒王后，此日青州土尚香。}
}
{
    \eA{[English source not present in the checked-in Bencraft/Joly files.]}
}
{
}

\Verse{209}
{
    \zA{众幕宾看了，便皆大赞：“小哥儿十三岁的人就如此，可知家学渊深真不诬矣。”}
}
{
    \eA{[English source not present in the checked-in Bencraft/Joly files.]}
}
{
}

\Verse{210}
{
    \zA{贾政笑道：“稚子口角，也还难为他。”}
}
{
    \eA{[English source not present in the checked-in Bencraft/Joly files.]}
}
{
}

\Verse{211}
{
    \zA{又看贾环的，是首五言律，写道是： 红粉不知愁，将军意未休。}
}
{
    \eA{[English source not present in the checked-in Bencraft/Joly files.]}
}
{
}

\Verse{212}
{
    \zA{掩啼离绣幕，抱恨出青州。}
}
{
    \eA{[English source not present in the checked-in Bencraft/Joly files.]}
}
{
}

\Verse{213}
{
    \zA{自谓酬王德，谁能复寇仇？}
}
{
    \eA{[English source not present in the checked-in Bencraft/Joly files.]}
}
{
}

\Verse{214}
{
    \zA{好题忠义墓，千古独风流。}
}
{
    \eA{[English source not present in the checked-in Bencraft/Joly files.]}
}
{
}

\Verse{215}
{
    \zA{众人道：“更佳。}
}
{
    \eA{[English source not present in the checked-in Bencraft/Joly files.]}
}
{
}

\Verse{216}
{
    \zA{到底大几岁年纪，立意又自不同。”}
}
{
    \eA{[English source not present in the checked-in Bencraft/Joly files.]}
}
{
}

\Verse{217}
{
    \zA{贾政道：“倒还不甚大错， 终不恳切。”}
}
{
    \eA{[English source not present in the checked-in Bencraft/Joly files.]}
}
{
}

\Verse{218}
{
    \zA{众人道：“这就罢了。}
}
{
    \eA{[English source not present in the checked-in Bencraft/Joly files.]}
}
{
}

\Verse{219}
{
    \zA{三爷才大不多几岁，俱在未冠之时。}
}
{
    \eA{[English source not present in the checked-in Bencraft/Joly files.]}
}
{
}

\Verse{220}
{
    \zA{如此用心 做去，再过几年，怕不是大阮小阮了么？”}
}
{
    \eA{[English source not present in the checked-in Bencraft/Joly files.]}
}
{
}

\Verse{221}
{
    \zA{贾政笑道：“过奖了。}
}
{
    \eA{[English source not present in the checked-in Bencraft/Joly files.]}
}
{
}

\Verse{222}
{
    \zA{只是不肯读书的 过失。”}
}
{
    \eA{[English source not present in the checked-in Bencraft/Joly files.]}
}
{
}

\Verse{223}
{
    \zA{因问宝玉。}
}
{
    \eA{[English source not present in the checked-in Bencraft/Joly files.]}
}
{
}

\Verse{224}
{
    \zA{众人道：“二爷细心镂刻，定又是风流悲感，不同此等的了。”}
}
{
    \eA{[English source not present in the checked-in Bencraft/Joly files.]}
}
{
}

\Verse{225}
{
    \zA{宝 玉笑道：“这个题目似不称近体，须得古体或歌或行长篇一首，方能恳切。”}
}
{
    \eA{[English source not present in the checked-in Bencraft/Joly files.]}
}
{
}

\Verse{226}
{
    \zA{众人 听了，都站起身来，点头拍手道：“我说他立意不同!每一题到手，必先度其体格 宜与不宜，这便是老手妙法。}
}
{
    \eA{[English source not present in the checked-in Bencraft/Joly files.]}
}
{
}

\Verse{227}
{
    \zA{这题目名曰《？？}
}
{
    \eA{[English source not present in the checked-in Bencraft/Joly files.]}
}
{
}

\Verse{228}
{
    \zA{词》，且既有了序，此必是长篇歌 行，方合体式。}
}
{
    \eA{[English source not present in the checked-in Bencraft/Joly files.]}
}
{
}

\Verse{229}
{
    \zA{或拟温八叉《击瓯歌》，或拟李长吉《会稽歌》，或拟白乐天《长 恨歌》，或拟咏古词，半叙半咏，流利飘逸，始能尽妙。”}
}
{
    \eA{[English source not present in the checked-in Bencraft/Joly files.]}
}
{
}

\Verse{230}
{
    \zA{贾政听说，也合了主意， 遂自提笔向纸上要写。}
}
{
    \eA{[English source not present in the checked-in Bencraft/Joly files.]}
}
{
}

\Verse{231}
{
    \zA{又向宝玉笑道：“如此甚好。}
}
{
    \eA{[English source not present in the checked-in Bencraft/Joly files.]}
}
{
}

\Verse{232}
{
    \zA{你念，我写。}
}
{
    \eA{[English source not present in the checked-in Bencraft/Joly files.]}
}
{
}

\Verse{233}
{
    \zA{若不好了，我捶 你的肉，谁许你先大言不惭的！”}
}
{
    \eA{[English source not present in the checked-in Bencraft/Joly files.]}
}
{
}

\Verse{234}
{
    \zA{宝玉只得念了一句道： 恒王好武兼好色， 贾政写了看时，摇头道：“粗鄙！”}
}
{
    \eA{[English source not present in the checked-in Bencraft/Joly files.]}
}
{
}

\Verse{235}
{
    \zA{一幕友道：“要这样方古，究竟不粗。}
}
{
    \eA{[English source not present in the checked-in Bencraft/Joly files.]}
}
{
}

\Verse{236}
{
    \zA{且看他 底下的。”}
}
{
    \eA{[English source not present in the checked-in Bencraft/Joly files.]}
}
{
}

\Verse{237}
{
    \zA{贾政道：“姑存之。”}
}
{
    \eA{[English source not present in the checked-in Bencraft/Joly files.]}
}
{
}

\Verse{238}
{
    \zA{宝玉又道： 遂教美女习骑射。}
}
{
    \eA{[English source not present in the checked-in Bencraft/Joly files.]}
}
{
}

\Verse{239}
{
    \zA{歌艳舞不成欢，列阵挽戈为自得。}
}
{
    \eA{[English source not present in the checked-in Bencraft/Joly files.]}
}
{
}

\Verse{240}
{
    \zA{贾政写出，众人都道：“只这第三句便古朴老健，极妙。}
}
{
    \eA{[English source not present in the checked-in Bencraft/Joly files.]}
}
{
}

\Verse{241}
{
    \zA{这第四句平叙，也最得体。”}
}
{
    \eA{[English source not present in the checked-in Bencraft/Joly files.]}
}
{
}

\Verse{242}
{
    \zA{贾政道：“休谬加奖誉，且看转的如何。”}
}
{
    \eA{[English source not present in the checked-in Bencraft/Joly files.]}
}
{
}

\Verse{243}
{
    \zA{宝玉念道： 眼前不见尘沙起，将军俏影红灯里。}
}
{
    \eA{[English source not present in the checked-in Bencraft/Joly files.]}
}
{
}

\Verse{244}
{
    \zA{众人听了这两句，便都叫妙：“好个‘不见尘沙起’!又承了一句‘俏影红灯里’， 用字用句皆入神化了。”}
}
{
    \eA{[English source not present in the checked-in Bencraft/Joly files.]}
}
{
}

\Verse{245}
{
    \zA{宝玉道： 叱咤时闻口舌香，霜矛雪剑娇难举。}
}
{
    \eA{[English source not present in the checked-in Bencraft/Joly files.]}
}
{
}

\Verse{246}
{
    \zA{众人听了更拍手笑道：“越发画出来了。}
}
{
    \eA{[English source not present in the checked-in Bencraft/Joly files.]}
}
{
}

\Verse{247}
{
    \zA{当日敢是宝公也在坐，见其娇而且闻其香？}
}
{
    \eA{[English source not present in the checked-in Bencraft/Joly files.]}
}
{
}

\Verse{248}
{
    \zA{不然何体贴至此。”}
}
{
    \eA{[English source not present in the checked-in Bencraft/Joly files.]}
}
{
}

\Verse{249}
{
    \zA{宝玉笑道：“闺阁习武，任其勇悍，怎似男人？}
}
{
    \eA{[English source not present in the checked-in Bencraft/Joly files.]}
}
{
}

\Verse{250}
{
    \zA{不问而可知娇 怯之形了。”}
}
{
    \eA{[English source not present in the checked-in Bencraft/Joly files.]}
}
{
}

\Verse{251}
{
    \zA{贾政道：“还不快续，这又有你说嘴的了？”}
}
{
    \eA{[English source not present in the checked-in Bencraft/Joly files.]}
}
{
}

\Verse{252}
{
    \zA{宝玉只得又想了一想， 念道： 丁香结子芙蓉绦， 众人都道：“转‘萧’韵更妙，这才流利飘逸。}
}
{
    \eA{[English source not present in the checked-in Bencraft/Joly files.]}
}
{
}

\Verse{253}
{
    \zA{而且这句子也绮靡秀媚得妙。”}
}
{
    \eA{[English source not present in the checked-in Bencraft/Joly files.]}
}
{
}

\Verse{254}
{
    \zA{贾 政写了，道：“这一句不好，已有过了‘口舌香’、‘娇难举’，何必又如此？}
}
{
    \eA{[English source not present in the checked-in Bencraft/Joly files.]}
}
{
}

\Verse{255}
{
    \zA{这 是力量不加，故又弄出这些堆砌货来搪塞。”}
}
{
    \eA{[English source not present in the checked-in Bencraft/Joly files.]}
}
{
}

\Verse{256}
{
    \zA{宝玉笑道：“长歌也须得要些词藻点 缀点缀，不然便觉萧索。”}
}
{
    \eA{[English source not present in the checked-in Bencraft/Joly files.]}
}
{
}

\Verse{257}
{
    \zA{贾政道：“你只顾说那些，这一句底下如何转至武事呢？}
}
{
    \eA{[English source not present in the checked-in Bencraft/Joly files.]}
}
{
}

\Verse{258}
{
    \zA{若再多说两句，岂不蛇足了？”}
}
{
    \eA{[English source not present in the checked-in Bencraft/Joly files.]}
}
{
}

\Verse{259}
{
    \zA{宝玉道：“如此，底下一句兜转煞住，想也使得。”}
}
{
    \eA{[English source not present in the checked-in Bencraft/Joly files.]}
}
{
}

\Verse{260}
{
    \zA{贾政冷笑道：“你有多大本领!上头说了一句大开门的散话，如今又要一句连转带 煞，岂不心有馀而力不足呢。”}
}
{
    \eA{[English source not present in the checked-in Bencraft/Joly files.]}
}
{
}

\Verse{261}
{
    \zA{宝玉听了，垂头想了一想，说了一句道： 不系明珠系宝刀。}
}
{
    \eA{[English source not present in the checked-in Bencraft/Joly files.]}
}
{
}

\Verse{262}
{
    \zA{忙问：“这一句可还使得？”}
}
{
    \eA{[English source not present in the checked-in Bencraft/Joly files.]}
}
{
}

\Verse{263}
{
    \zA{众人拍案叫绝。}
}
{
    \eA{[English source not present in the checked-in Bencraft/Joly files.]}
}
{
}

\Verse{264}
{
    \zA{贾政笑道：“且放着，再续。”}
}
{
    \eA{[English source not present in the checked-in Bencraft/Joly files.]}
}
{
}

\Verse{265}
{
    \zA{宝玉 道：“使得，我便一气连下去了；若使不得，索性涂了，我再想别的意思出来，再 另措词。”}
}
{
    \eA{[English source not present in the checked-in Bencraft/Joly files.]}
}
{
}

\Verse{266}
{
    \zA{贾政听了，便喝道：“多话!不好了再做。}
}
{
    \eA{[English source not present in the checked-in Bencraft/Joly files.]}
}
{
}

\Verse{267}
{
    \zA{便做十篇百篇，还怕辛苦了 不成？”}
}
{
    \eA{[English source not present in the checked-in Bencraft/Joly files.]}
}
{
}

\Verse{268}
{
    \zA{宝玉听了，只得想了一会，便念道： 战罢夜阑心力怯，脂痕粉渍污鲛绡。}
}
{
    \eA{[English source not present in the checked-in Bencraft/Joly files.]}
}
{
}

\Verse{269}
{
    \zA{贾政道：“这又是一段了。}
}
{
    \eA{[English source not present in the checked-in Bencraft/Joly files.]}
}
{
}

\Verse{270}
{
    \zA{底下怎么样？”}
}
{
    \eA{[English source not present in the checked-in Bencraft/Joly files.]}
}
{
}

\Verse{271}
{
    \zA{宝玉道： 明年流寇走山东，强吞虎豹势如蜂。}
}
{
    \eA{[English source not present in the checked-in Bencraft/Joly files.]}
}
{
}

\Verse{272}
{
    \zA{众人道：“好个‘走’字，便见得高低了。}
}
{
    \eA{[English source not present in the checked-in Bencraft/Joly files.]}
}
{
}

\Verse{273}
{
    \zA{且通句转的也不板。”}
}
{
    \eA{[English source not present in the checked-in Bencraft/Joly files.]}
}
{
}

\Verse{274}
{
    \zA{宝玉又念道： 王率天兵思剿灭，一战再战不成功。}
}
{
    \eA{[English source not present in the checked-in Bencraft/Joly files.]}
}
{
}

\Verse{275}
{
    \zA{腥风吹折陇中麦，日照旌旗虎帐空。}
}
{
    \eA{[English source not present in the checked-in Bencraft/Joly files.]}
}
{
}

\Verse{276}
{
    \zA{青山寂寂水澌澌，正是恒王战死时。}
}
{
    \eA{[English source not present in the checked-in Bencraft/Joly files.]}
}
{
}

\Verse{277}
{
    \zA{雨淋白骨血染草，月冷黄昏鬼守尸。}
}
{
    \eA{[English source not present in the checked-in Bencraft/Joly files.]}
}
{
}

\Verse{278}
{
    \zA{众人都道：“妙极，妙极!布置叙事词藻，无不尽美。}
}
{
    \eA{[English source not present in the checked-in Bencraft/Joly files.]}
}
{
}

\Verse{279}
{
    \zA{且看如何至四娘，必另有妙 转奇句。”}
}
{
    \eA{[English source not present in the checked-in Bencraft/Joly files.]}
}
{
}

\Verse{280}
{
    \zA{宝玉又念道： 纷纷将士只保身，青州眼见皆灰尘。}
}
{
    \eA{[English source not present in the checked-in Bencraft/Joly files.]}
}
{
}

\Verse{281}
{
    \zA{不期忠义明闺阁，愤起恒王得意人。}
}
{
    \eA{[English source not present in the checked-in Bencraft/Joly files.]}
}
{
}

\Verse{282}
{
    \zA{众人都道：“铺叙得委婉！”}
}
{
    \eA{[English source not present in the checked-in Bencraft/Joly files.]}
}
{
}

\Verse{283}
{
    \zA{贾政道：“太多了，底下只怕累赘呢。”}
}
{
    \eA{[English source not present in the checked-in Bencraft/Joly files.]}
}
{
}

\Verse{284}
{
    \zA{宝玉又道： 恒王得意数谁行：？？}
}
{
    \eA{[English source not present in the checked-in Bencraft/Joly files.]}
}
{
}

\Verse{285}
{
    \zA{将军林四娘。}
}
{
    \eA{[English source not present in the checked-in Bencraft/Joly files.]}
}
{
}

\Verse{286}
{
    \zA{号令秦姬驱赵女，？}
}
{
    \eA{[English source not present in the checked-in Bencraft/Joly files.]}
}
{
}

\Verse{287}
{
    \zA{桃艳李临疆场。}
}
{
    \eA{[English source not present in the checked-in Bencraft/Joly files.]}
}
{
}

\Verse{288}
{
    \zA{绣鞍有泪春愁重，铁甲无声夜气凉。}
}
{
    \eA{[English source not present in the checked-in Bencraft/Joly files.]}
}
{
}

\Verse{289}
{
    \zA{胜负自难先预定，誓盟生死报前王。}
}
{
    \eA{[English source not present in the checked-in Bencraft/Joly files.]}
}
{
}

\Verse{290}
{
    \zA{贼势猖獗不可敌，柳折花残血凝碧。}
}
{
    \eA{[English source not present in the checked-in Bencraft/Joly files.]}
}
{
}

\Verse{291}
{
    \zA{马践胭脂骨髓香，魂依城郭家乡隔。}
}
{
    \eA{[English source not present in the checked-in Bencraft/Joly files.]}
}
{
}

\Verse{292}
{
    \zA{星驰时报入京师，谁家儿女不伤悲! 天子惊慌愁失守，此时文武皆垂首。}
}
{
    \eA{[English source not present in the checked-in Bencraft/Joly files.]}
}
{
}

\Verse{293}
{
    \zA{何事文武立朝纲，不及闺中林四娘？}
}
{
    \eA{[English source not present in the checked-in Bencraft/Joly files.]}
}
{
}

\Verse{294}
{
    \zA{我为四娘长叹息，歌成馀意尚彷徨! 念毕，众人都大赞不止。}
}
{
    \eA{[English source not present in the checked-in Bencraft/Joly files.]}
}
{
}

\Verse{295}
{
    \zA{又从头看了一遍。}
}
{
    \eA{[English source not present in the checked-in Bencraft/Joly files.]}
}
{
}

\Verse{296}
{
    \zA{贾政笑道：“虽说了几句，到底不大恳 切。”}
}
{
    \eA{[English source not present in the checked-in Bencraft/Joly files.]}
}
{
}

\Verse{297}
{
    \zA{因说：“去罢。”}
}
{
    \eA{[English source not present in the checked-in Bencraft/Joly files.]}
}
{
}

\Verse{298}
{
    \zA{三人如放了赦的一般，一齐出来，各自回房。}
}
{
    \eA{[English source not present in the checked-in Bencraft/Joly files.]}
}
{
}

\Verse{299}
{
    \zA{众人皆无别 话，不过至晚安歇而已。}
}
{
    \eA{[English source not present in the checked-in Bencraft/Joly files.]}
}
{
}

\Verse{300}
{
    \zA{独有宝玉，一心凄楚。}
}
{
    \eA{[English source not present in the checked-in Bencraft/Joly files.]}
}
{
}

\Verse{301}
{
    \zA{回至园中，猛见池上芙蓉，想起小丫鬟说晴雯做了芙蓉 之神，不觉又喜欢起来，乃看着芙蓉嗟叹了一会。}
}
{
    \eA{[English source not present in the checked-in Bencraft/Joly files.]}
}
{
}

\Verse{302}
{
    \zA{忽又想起：“死后并未至灵前一 祭，如今何不在芙蓉前一祭，岂不尽了礼？”}
}
{
    \eA{[English source not present in the checked-in Bencraft/Joly files.]}
}
{
}

\Verse{303}
{
    \zA{想毕，便欲行礼。}
}
{
    \eA{[English source not present in the checked-in Bencraft/Joly files.]}
}
{
}

\Verse{304}
{
    \zA{忽又止道：“虽如 此，亦不可太草率了，须得衣冠整齐，奠仪周备，方为诚敬。”}
}
{
    \eA{[English source not present in the checked-in Bencraft/Joly files.]}
}
{
}

\Verse{305}
{
    \zA{想了一想：“古人 云，‘潢污行潦，荇藻苹蘩之贱，可以羞王公，荐鬼神’，原不在物之贵贱，只在 心之诚敬而已。}
}
{
    \eA{[English source not present in the checked-in Bencraft/Joly files.]}
}
{
}

\Verse{306}
{
    \zA{然非自作一篇诔文，这一段凄惨酸楚，竟无处可以发泄了。”}
}
{
    \eA{[English source not present in the checked-in Bencraft/Joly files.]}
}
{
}

\Verse{307}
{
    \zA{因用 晴雯素日所喜之冰鲛？}
}
{
    \eA{[English source not present in the checked-in Bencraft/Joly files.]}
}
{
}

\Verse{308}
{
    \zA{一幅，楷字写成，名曰《芙蓉女儿诔》，前序后歌；又备了 晴雯素喜的四样吃食。}
}
{
    \eA{[English source not present in the checked-in Bencraft/Joly files.]}
}
{
}

\Verse{309}
{
    \zA{于是黄昏人静之时，命那小丫头捧至芙蓉前，先行礼毕，将 那诔文即挂于芙蓉枝上，乃泣涕念曰：
        维太平不易之元，蓉桂竞芳之月，无可奈何之日，怡红院浊玉谨以群花之蕊、 冰鲛之？}
}
{
    \eA{[English source not present in the checked-in Bencraft/Joly files.]}
}
{
}

\Verse{310}
{
    \zA{、沁芳之泉、枫露之茗：四者虽微，聊以达诚申信，乃致祭于白帝宫中抚 司秋艳芙蓉女儿之前曰： 窃思女儿自临人世，迄今凡十有六载。}
}
{
    \eA{[English source not present in the checked-in Bencraft/Joly files.]}
}
{
}

\Verse{311}
{
    \zA{其先之乡籍姓氏，湮沦而莫能考者久矣。}
}
{
    \eA{[English source not present in the checked-in Bencraft/Joly files.]}
}
{
}

\Verse{312}
{
    \zA{而玉得于衾枕栉沐之间，栖息宴游之夕，亲昵狎亵，相与共处者，仅五年八月有奇。}
}
{
    \eA{[English source not present in the checked-in Bencraft/Joly files.]}
}
{
}

\Verse{313}
{
    \zA{忆女曩生之昔，其为质则金玉不足喻其贵，其为体则冰雪不足喻其洁。}
}
{
    \eA{[English source not present in the checked-in Bencraft/Joly files.]}
}
{
}

\Verse{314}
{
    \zA{其为神则星 日不足喻其精，其为貌则花月不足喻其色。}
}
{
    \eA{[English source not present in the checked-in Bencraft/Joly files.]}
}
{
}

\Verse{315}
{
    \zA{姊娣悉慕？}
}
{
    \eA{[English source not present in the checked-in Bencraft/Joly files.]}
}
{
}

\Verse{316}
{
    \zA{娴，妪媪咸仰慧德。}
}
{
    \eA{[English source not present in the checked-in Bencraft/Joly files.]}
}
{
}

\Verse{317}
{
    \zA{孰料鸠 鸩恶其高，鹰鸷翻遭？？}
}
{
    \eA{[English source not present in the checked-in Bencraft/Joly files.]}
}
{
}

\Verse{318}
{
    \zA{；？？}
}
{
    \eA{[English source not present in the checked-in Bencraft/Joly files.]}
}
{
}

\Verse{319}
{
    \zA{妒其臭，？}
}
{
    \eA{[English source not present in the checked-in Bencraft/Joly files.]}
}
{
}

\Verse{320}
{
    \zA{兰竟被芟？。}
}
{
    \eA{[English source not present in the checked-in Bencraft/Joly files.]}
}
{
}

\Verse{321}
{
    \zA{花原自怯，岂奈狂飚？}
}
{
    \eA{[English source not present in the checked-in Bencraft/Joly files.]}
}
{
}

\Verse{322}
{
    \zA{柳本 多愁，何禁骤雨!偶遭蛊虿之谗，遂抱膏肓之疾。}
}
{
    \eA{[English source not present in the checked-in Bencraft/Joly files.]}
}
{
}

\Verse{323}
{
    \zA{故樱唇红褪，韵吐呻吟；杏脸香枯， 色陈？}
}
{
    \eA{[English source not present in the checked-in Bencraft/Joly files.]}
}
{
}

\Verse{324}
{
    \zA{颔。}
}
{
    \eA{[English source not present in the checked-in Bencraft/Joly files.]}
}
{
}

\Verse{325}
{
    \zA{诼谣？}
}
{
    \eA{[English source not present in the checked-in Bencraft/Joly files.]}
}
{
}

\Verse{326}
{
    \zA{诟，出自屏帷；荆棘蓬榛，蔓延窗户。}
}
{
    \eA{[English source not present in the checked-in Bencraft/Joly files.]}
}
{
}

\Verse{327}
{
    \zA{既怀幽沉于不尽，复含罔屈 于无穷。}
}
{
    \eA{[English source not present in the checked-in Bencraft/Joly files.]}
}
{
}

\Verse{328}
{
    \zA{高标见嫉，闺闱恨比长沙；贞烈遭危，巾帼惨于雁塞。}
}
{
    \eA{[English source not present in the checked-in Bencraft/Joly files.]}
}
{
}

\Verse{329}
{
    \zA{自蓄辛酸，谁怜夭折？}
}
{
    \eA{[English source not present in the checked-in Bencraft/Joly files.]}
}
{
}

\Verse{330}
{
    \zA{仙云既散，芳趾难寻。}
}
{
    \eA{[English source not present in the checked-in Bencraft/Joly files.]}
}
{
}

\Verse{331}
{
    \zA{洲迷聚窟，何来却死之香？}
}
{
    \eA{[English source not present in the checked-in Bencraft/Joly files.]}
}
{
}

\Verse{332}
{
    \zA{海失灵槎，不获回生之药。}
}
{
    \eA{[English source not present in the checked-in Bencraft/Joly files.]}
}
{
}

\Verse{333}
{
    \zA{眉黛烟 青，昨犹我画；指环玉冷，今倩谁温？}
}
{
    \eA{[English source not present in the checked-in Bencraft/Joly files.]}
}
{
}

\Verse{334}
{
    \zA{鼎炉之剩药犹存，襟泪之馀痕尚渍。}
}
{
    \eA{[English source not present in the checked-in Bencraft/Joly files.]}
}
{
}

\Verse{335}
{
    \zA{镜分鸾影， 愁开麝月之奁；梳化龙飞，哀折檀云之齿。}
}
{
    \eA{[English source not present in the checked-in Bencraft/Joly files.]}
}
{
}

\Verse{336}
{
    \zA{委金钿于草莽，拾翠盒于尘埃。}
}
{
    \eA{[English source not present in the checked-in Bencraft/Joly files.]}
}
{
}

\Verse{337}
{
    \zA{楼空？？}
}
{
    \eA{[English source not present in the checked-in Bencraft/Joly files.]}
}
{
}

\Verse{338}
{
    \zA{鹊，从悬七夕之针；带断鸳鸯，谁续五丝之缕？}
}
{
    \eA{[English source not present in the checked-in Bencraft/Joly files.]}
}
{
}

\Verse{339}
{
    \zA{况乃金天属节，白帝司时；孤衾有 梦，空室无人。}
}
{
    \eA{[English source not present in the checked-in Bencraft/Joly files.]}
}
{
}

\Verse{340}
{
    \zA{桐阶月暗，芳魂与倩影同消；蓉帐香残，娇喘共细腰俱绝。}
}
{
    \eA{[English source not present in the checked-in Bencraft/Joly files.]}
}
{
}

\Verse{341}
{
    \zA{连天衰 草，岂独蒹葭；匝地悲声，无非蟋蟀。}
}
{
    \eA{[English source not present in the checked-in Bencraft/Joly files.]}
}
{
}

\Verse{342}
{
    \zA{露阶晚砌，穿帘不度寒砧；雨荔秋垣，隔院 希闻怨笛。}
}
{
    \eA{[English source not present in the checked-in Bencraft/Joly files.]}
}
{
}

\Verse{343}
{
    \zA{芳名未泯，檐前鹦鹉犹呼；艳质将亡，槛外海棠预萎。}
}
{
    \eA{[English source not present in the checked-in Bencraft/Joly files.]}
}
{
}

\Verse{344}
{
    \zA{捉迷屏后，莲瓣 无声；斗草庭前，兰芳枉待。}
}
{
    \eA{[English source not present in the checked-in Bencraft/Joly files.]}
}
{
}

\Verse{345}
{
    \zA{抛残绣线，银笺彩袖谁裁？}
}
{
    \eA{[English source not present in the checked-in Bencraft/Joly files.]}
}
{
}

\Verse{346}
{
    \zA{折断冰丝，金斗御香未熨。}
}
{
    \eA{[English source not present in the checked-in Bencraft/Joly files.]}
}
{
}

\Verse{347}
{
    \zA{昨承严命，既趋车而远陟芳园；今犯慈威，复拄杖而遣抛孤柩。}
}
{
    \eA{[English source not present in the checked-in Bencraft/Joly files.]}
}
{
}

\Verse{348}
{
    \zA{及闻蕙棺被燹，顿 违共穴之情；石椁成灾，愧逮同灰之诮。}
}
{
    \eA{[English source not present in the checked-in Bencraft/Joly files.]}
}
{
}

\Verse{349}
{
    \zA{尔乃西风古寺，淹滞青磷；落日荒丘，零 星白骨。}
}
{
    \eA{[English source not present in the checked-in Bencraft/Joly files.]}
}
{
}

\Verse{350}
{
    \zA{楸榆飒飒，蓬艾萧萧。}
}
{
    \eA{[English source not present in the checked-in Bencraft/Joly files.]}
}
{
}

\Verse{351}
{
    \zA{隔雾圹以啼猿，绕烟塍而泣鬼。}
}
{
    \eA{[English source not present in the checked-in Bencraft/Joly files.]}
}
{
}

\Verse{352}
{
    \zA{岂道红绡帐里，公 子情深；始信黄土陇中，女儿命薄!汝南斑斑泪血，洒向西风；梓泽默默馀衷，诉 凭冷月。}
}
{
    \eA{[English source not present in the checked-in Bencraft/Joly files.]}
}
{
}

\Verse{353}
{
    \zA{呜呼!固鬼蜮之为灾，岂神灵之有妒!毁？}
}
{
    \eA{[English source not present in the checked-in Bencraft/Joly files.]}
}
{
}

\Verse{354}
{
    \zA{奴之口，讨岂从宽？}
}
{
    \eA{[English source not present in the checked-in Bencraft/Joly files.]}
}
{
}

\Verse{355}
{
    \zA{剖悍妇之心， 忿犹未释。}
}
{
    \eA{[English source not present in the checked-in Bencraft/Joly files.]}
}
{
}

\Verse{356}
{
    \zA{在卿之尘缘虽浅，而玉之鄙意尤深。}
}
{
    \eA{[English source not present in the checked-in Bencraft/Joly files.]}
}
{
}

\Verse{357}
{
    \zA{因蓄？？}
}
{
    \eA{[English source not present in the checked-in Bencraft/Joly files.]}
}
{
}

\Verse{358}
{
    \zA{之思，不禁谆谆之问。}
}
{
    \eA{[English source not present in the checked-in Bencraft/Joly files.]}
}
{
}

\Verse{359}
{
    \zA{始 知上帝垂旌，花宫待诏。}
}
{
    \eA{[English source not present in the checked-in Bencraft/Joly files.]}
}
{
}

\Verse{360}
{
    \zA{生侪兰蕙，死辖芙蓉。}
}
{
    \eA{[English source not present in the checked-in Bencraft/Joly files.]}
}
{
}

\Verse{361}
{
    \zA{听小婢之言，似涉无稽；据浊玉之思， 深为有据。}
}
{
    \eA{[English source not present in the checked-in Bencraft/Joly files.]}
}
{
}

\Verse{362}
{
    \zA{何也：昔叶法善摄魂以撰碑，李长吉被诏而为记：事虽殊，其理则一也。}
}
{
    \eA{[English source not present in the checked-in Bencraft/Joly files.]}
}
{
}

\Verse{363}
{
    \zA{故相物以配才，苟非其人，恶乃滥乎？}
}
{
    \eA{[English source not present in the checked-in Bencraft/Joly files.]}
}
{
}

\Verse{364}
{
    \zA{始信上帝委托权衡，可谓至洽至协，庶不负其 所秉赋也。}
}
{
    \eA{[English source not present in the checked-in Bencraft/Joly files.]}
}
{
}

\Verse{365}
{
    \zA{因希其不昧之灵，或陟降于兹，特不揣鄙俗之词，有污慧听。}
}
{
    \eA{[English source not present in the checked-in Bencraft/Joly files.]}
}
{
}

\Verse{366}
{
    \zA{乃歌而招 之曰： 天何如是之苍苍兮，乘玉虬以游乎穹窿耶？}
}
{
    \eA{[English source not present in the checked-in Bencraft/Joly files.]}
}
{
}

\Verse{367}
{
    \zA{地何如是之茫茫兮，驾瑶象以降乎 泉壤耶？}
}
{
    \eA{[English source not present in the checked-in Bencraft/Joly files.]}
}
{
}

\Verse{368}
{
    \zA{望伞盖之陆离兮，抑箕尾之光耶？}
}
{
    \eA{[English source not present in the checked-in Bencraft/Joly files.]}
}
{
}

\Verse{369}
{
    \zA{列羽葆而为前导兮，卫危虚于傍耶？}
}
{
    \eA{[English source not present in the checked-in Bencraft/Joly files.]}
}
{
}

\Verse{370}
{
    \zA{驱丰 隆以为庇从兮，望舒月以临耶？}
}
{
    \eA{[English source not present in the checked-in Bencraft/Joly files.]}
}
{
}

\Verse{371}
{
    \zA{听车轨而伊轧兮：御鸾？}
}
{
    \eA{[English source not present in the checked-in Bencraft/Joly files.]}
}
{
}

\Verse{372}
{
    \zA{以征耶？}
}
{
    \eA{[English source not present in the checked-in Bencraft/Joly files.]}
}
{
}

\Verse{373}
{
    \zA{闻馥郁而飘然兮， 纫蘅杜以为佩耶？}
}
{
    \eA{[English source not present in the checked-in Bencraft/Joly files.]}
}
{
}

\Verse{374}
{
    \zA{斓裙裾之烁烁兮，镂明月以为？}
}
{
    \eA{[English source not present in the checked-in Bencraft/Joly files.]}
}
{
}

\Verse{375}
{
    \zA{耶？}
}
{
    \eA{[English source not present in the checked-in Bencraft/Joly files.]}
}
{
}

\Verse{376}
{
    \zA{借葳蕤而成坛？}
}
{
    \eA{[English source not present in the checked-in Bencraft/Joly files.]}
}
{
}

\Verse{377}
{
    \zA{兮，檠莲焰以 烛兰膏耶？}
}
{
    \eA{[English source not present in the checked-in Bencraft/Joly files.]}
}
{
}

\Verse{378}
{
    \zA{文瓠匏以为觯？}
}
{
    \eA{[English source not present in the checked-in Bencraft/Joly files.]}
}
{
}

\Verse{379}
{
    \zA{兮，洒？？}
}
{
    \eA{[English source not present in the checked-in Bencraft/Joly files.]}
}
{
}

\Verse{380}
{
    \zA{以浮桂醑耶？}
}
{
    \eA{[English source not present in the checked-in Bencraft/Joly files.]}
}
{
}

\Verse{381}
{
    \zA{瞻云气而凝眸兮，仿佛有所觇耶？}
}
{
    \eA{[English source not present in the checked-in Bencraft/Joly files.]}
}
{
}

\Verse{382}
{
    \zA{俯波痕而属耳兮，恍惚有所闻耶？}
}
{
    \eA{[English source not present in the checked-in Bencraft/Joly files.]}
}
{
}

\Verse{383}
{
    \zA{期汗漫而无际兮，捐弃予于尘埃耶？}
}
{
    \eA{[English source not present in the checked-in Bencraft/Joly files.]}
}
{
}

\Verse{384}
{
    \zA{倩风廉之为余 驱车兮，冀联辔而携归耶？}
}
{
    \eA{[English source not present in the checked-in Bencraft/Joly files.]}
}
{
}

\Verse{385}
{
    \zA{余中心为之慨然兮，徒？？}
}
{
    \eA{[English source not present in the checked-in Bencraft/Joly files.]}
}
{
}

\Verse{386}
{
    \zA{而何为耶？}
}
{
    \eA{[English source not present in the checked-in Bencraft/Joly files.]}
}
{
}

\Verse{387}
{
    \zA{卿偃然而长寝兮， 岂天运之变于斯耶？}
}
{
    \eA{[English source not present in the checked-in Bencraft/Joly files.]}
}
{
}

\Verse{388}
{
    \zA{既窀穸且安稳兮，反其真而又奚化耶？}
}
{
    \eA{[English source not present in the checked-in Bencraft/Joly files.]}
}
{
}

\Verse{389}
{
    \zA{余犹桎梏而悬附兮，灵格 余以嗟来耶？}
}
{
    \eA{[English source not present in the checked-in Bencraft/Joly files.]}
}
{
}

\Verse{390}
{
    \zA{来兮止兮，卿其来耶？}
}
{
    \eA{[English source not present in the checked-in Bencraft/Joly files.]}
}
{
}

\Verse{391}
{
    \zA{若夫鸿蒙而居，寂静以处，虽临于兹，余亦莫睹。}
}
{
    \eA{[English source not present in the checked-in Bencraft/Joly files.]}
}
{
}

\Verse{392}
{
    \zA{搴烟萝而为步障，列苍蒲而 森行伍。}
}
{
    \eA{[English source not present in the checked-in Bencraft/Joly files.]}
}
{
}

\Verse{393}
{
    \zA{警柳眼之贪眠，释莲心之味苦，素女约于桂岩，宓妃迎于兰渚。}
}
{
    \eA{[English source not present in the checked-in Bencraft/Joly files.]}
}
{
}

\Verse{394}
{
    \zA{弄玉吹笙， 寒簧击？。}
}
{
    \eA{[English source not present in the checked-in Bencraft/Joly files.]}
}
{
}

\Verse{395}
{
    \zA{征嵩岳之妃，启骊山之姥。}
}
{
    \eA{[English source not present in the checked-in Bencraft/Joly files.]}
}
{
}

\Verse{396}
{
    \zA{龟呈洛浦之灵，兽作咸池之舞。}
}
{
    \eA{[English source not present in the checked-in Bencraft/Joly files.]}
}
{
}

\Verse{397}
{
    \zA{潜赤水兮龙 吟，集珠林兮凤翥。}
}
{
    \eA{[English source not present in the checked-in Bencraft/Joly files.]}
}
{
}

\Verse{398}
{
    \zA{爰格爰诚，匪笤匪？。}
}
{
    \eA{[English source not present in the checked-in Bencraft/Joly files.]}
}
{
}

\Verse{399}
{
    \zA{发轫乎霞城，还旌乎玄圃。}
}
{
    \eA{[English source not present in the checked-in Bencraft/Joly files.]}
}
{
}

\Verse{400}
{
    \zA{既显微而若 逋，复氤氲而倏阻。}
}
{
    \eA{[English source not present in the checked-in Bencraft/Joly files.]}
}
{
}

\Verse{401}
{
    \zA{离合兮烟云，空蒙兮雾雨。}
}
{
    \eA{[English source not present in the checked-in Bencraft/Joly files.]}
}
{
}

\Verse{402}
{
    \zA{尘霾敛兮星高，溪山丽兮月午。}
}
{
    \eA{[English source not present in the checked-in Bencraft/Joly files.]}
}
{
}

\Verse{403}
{
    \zA{何 心意之怦怦，若寤寐之栩栩？}
}
{
    \eA{[English source not present in the checked-in Bencraft/Joly files.]}
}
{
}

\Verse{404}
{
    \zA{余乃欷？}
}
{
    \eA{[English source not present in the checked-in Bencraft/Joly files.]}
}
{
}

\Verse{405}
{
    \zA{怅怏，泣涕彷徨。}
}
{
    \eA{[English source not present in the checked-in Bencraft/Joly files.]}
}
{
}

\Verse{406}
{
    \zA{人语兮寂历，天籁兮？？。}
}
{
    \eA{[English source not present in the checked-in Bencraft/Joly files.]}
}
{
}

\Verse{407}
{
    \zA{鸟惊散而飞，鱼唼喋以响。}
}
{
    \eA{[English source not present in the checked-in Bencraft/Joly files.]}
}
{
}

\Verse{408}
{
    \zA{志哀兮是祷，成礼兮期祥。}
}
{
    \eA{[English source not present in the checked-in Bencraft/Joly files.]}
}
{
}

\Verse{409}
{
    \zA{呜呼哀哉!尚飨! 读毕，遂焚帛奠茗，依依不舍。}
}
{
    \eA{[English source not present in the checked-in Bencraft/Joly files.]}
}
{
}

\Verse{410}
{
    \zA{小丫鬟催至再四，方才回身。}
}
{
    \eA{[English source not present in the checked-in Bencraft/Joly files.]}
}
{
}

\Verse{411}
{
    \zA{忽听山石之后有一人笑道：“且请留步。”}
}
{
    \eA{[English source not present in the checked-in Bencraft/Joly files.]}
}
{
}

\Verse{412}
{
    \zA{二人听了，不觉大惊。}
}
{
    \eA{[English source not present in the checked-in Bencraft/Joly files.]}
}
{
}

\Verse{413}
{
    \zA{那小丫鬟回 头一看，却是个人影儿从芙蓉花里走出来，他便大叫：“不好，有鬼!晴雯真来显 魂了！”}
}
{
    \eA{[English source not present in the checked-in Bencraft/Joly files.]}
}
{
}

\Verse{414}
{
    \zA{唬得宝玉也忙看时，―― 究竟是人是鬼，下回分解。}
}
{
    \eA{[English source not present in the checked-in Bencraft/Joly files.]}
}
{
}

\Verse{415}
{
    \zA{(本章完)}
}
{
    \eA{[English source not present in the checked-in Bencraft/Joly files.]}
}
{
}
